
\subsubsection{6.1 Interpretation of Findings}

The results establish three empirical patterns. First, human follow-up turns in HH-RLHF helpfulness conversations exhibit substantially higher semantic similarity to their actual AI partner turns than to random or topic-matched baseline turns. The effect size relative to the random baseline is large (d = 1.998), and the effect remains present when compared to KNN topic-matched controls (d = 0.400--0.417 across models). This indicates that the measured similarity is not fully attributable to topical co-occurrence, though SBERT embeddings do not permit clean separation of topical and stylistic components.

Second, C\_sem increases monotonically from T1 to T3 across all three SBERT models, with J-T p = 0.001 in each case. The effect size for the T1-to-T3 contrast is small to medium (d = 0.183--0.254). Adjacent-tier contrasts are smaller (d = 0.087--0.114 for T1 vs. T2), and some fall below the pre-specified 0.1 threshold. IPW correction for distributional shift across tiers preserves the monotonic ordering.

Third, human-to-AI accommodation exceeds AI-to-human accommodation in all nine tier-by-model conditions. The effect sizes are variable (d = 0.061--0.405) and tend to decrease for higher-quality tiers.

The two falsified mechanism hypotheses provide constraints on interpretation. The within-prompt quality probe (h-m3) yields a reversed signal: human follow-up turns are more similar to the rejected AI response than to the chosen one. This reversal is strongest in the highest-quality tier (T3, d up to -0.738). One interpretation is that rejected responses in HH-RLHF tend to be longer and more topically expansive, providing more semantic surface area that incidentally overlaps with the human's subsequent turn, independent of the quality dimensions on which RLHF raters based their preference annotations. The politeness-style mediation analysis (h-m4) finds no evidence that a cosine-based politeness proxy explains C\_sem variance.

These falsifications indicate that the tier-accommodation association operates at the population level -- conversations drawn from higher-quality tiers exhibit stronger accommodation on average -- rather than through a within-conversation mechanism where humans perceive and respond to individual response quality.

\subsubsection{6.2 Limitations}

\textbf{Cross-sectional design.} The HH-RLHF dataset lacks user identifiers; each conversation is anonymous. The finding that higher-tier conversations show stronger C\_sem cannot distinguish accommodation from user self-selection: more sophisticated users may self-select into higher-tier conversations and independently exhibit communication patterns that yield higher C\_sem. Within-user longitudinal data (e.g., from LMSYS Chatbot Arena with session identifiers) would be needed to test within-user accommodation trajectories.

\textbf{Topical-stylistic entanglement in SBERT.} SBERT embeddings capture both content and style. C\_sem cannot be decomposed into topical and stylistic components. The KNN topic-matched baseline provides a partial control (d = 0.417 above topic-matched), but residual topical similarity may contribute to the measured effect. Style-factored sentence representations would enable more precise decomposition.

\textbf{Tier confounds.} The three HH-RLHF splits were collected at different stages of Anthropic's deployment pipeline. Differences in user demographics, conversation topics, and interaction conventions across tiers may confound the RLHF quality variable. IPW correction addresses measured distributional shift but cannot control for unmeasured confounds.

\textbf{Unresolved mechanism.} Both tested mechanism hypotheses (within-prompt quality discrimination, politeness-style mediation) are falsified. The proximal mechanism linking RLHF tier to population-level accommodation remains unidentified. This limits the theoretical interpretation of the tier-scaling finding.

\textbf{PM-proxy operationalization.} The politeness-style proxy used in h-m4 is based on cosine similarity to a hand-curated centroid, which may not capture the relevant dimensions of AI response quality. However, the essentially zero coefficient (beta < 1e-4) and very low model R-squared (less than 0.012) suggest that even a better proxy would need to explain substantially more variance than the current feature set.

\subsubsection{6.3 Implications}

The finding that RLHF tier quality co-varies with human semantic behavior at the population level suggests that RLHF evaluation may benefit from incorporating measures of human-side communicative patterns, in addition to standard AI output quality metrics. The h-m3 reversal (human follow-ups are more similar to rejected than chosen responses) indicates that RLHF's chosen/rejected annotation captures quality dimensions that are partly orthogonal to conversational semantic continuity. Researchers using the HH-RLHF chosen/rejected structure as a proxy for conversational effectiveness should be aware of this dissociation.
